\documentclass[abstract=on,10pt,a4paper,bibliography=totocnumbered]{article}
\usepackage[paper=a4paper,left=35mm,right=35mm,top=25mm,bottom=30mm]{geometry}
\usepackage[doublespacing]{setspace}
\usepackage[english]{babel}
\usepackage[utf8]{inputenc}
\usepackage{amsmath}
\usepackage{colortbl}
\usepackage{amsfonts}
\usepackage{amssymb}
\usepackage{gensymb}
\usepackage{textcomp}
\usepackage{graphicx}
\usepackage{tikz}
\usepackage{enumerate}
\usepackage{enumitem}
\usepackage{subcaption}
\usepackage{booktabs}
\usepackage[hidelinks]{hyperref}
\usepackage[nameinlink]{cleveref}
\usepackage{multirow}
\usepackage{arydshln}
\usepackage[flushleft]{threeparttable}
% \usepackage[nomarkers, nolists]{endfloat}
\usepackage{scalerel}
\usepackage{makecell}
\usepackage{ifthen}
\usepackage{tikz}
\usepackage{lineno}
\usepackage{adjustbox}
\usetikzlibrary{svg.path}

%------------------------------------------------------------------------------
%  Bibtex vs Biblatex
%------------------------------------------------------------------------------
% While bibtex is very nice to use, it's insanely slow. I therefore work with
% bibtex until the final stage. Set the value below to "true" or "false" to pick
% which one should be used.
\newboolean{usebiblatex}
\setboolean{usebiblatex}{true}

% Depending on the question, the commands below will prepare the document
% accordingly
\ifthenelse{\boolean{usebiblatex}}{
  %-----------------------------------------------------------------------------
  %  If using biblatex
  %-----------------------------------------------------------------------------
  % Load required packages and refer to correct bibliography
  \usepackage{csquotes}
  \usepackage[
      backend=biber
    , style=apa          % To display Author-year in the text
    , natbib=true        % Necessary to use citep, etc.
    , hyperref=true      % To have clickable links in the references
    , useeditor=false    % Don't add editor to publications
    , sorting=ynt        % Sort by year, name, and title
    , uniquename=false   % Avoid that initials are put
  ]{biblatex}
  \addbibresource{Literature.bib}

  % To include author in link
  \DeclareFieldFormat{citehyperref}{%
  \DeclareFieldAlias{bibhyperref}{noformat}% Avoid nested links
  \bibhyperref{#1}}
  \DeclareFieldFormat{textcitehyperref}{%
  \DeclareFieldAlias{bibhyperref}{noformat}% Avoid nested links
  \bibhyperref{%
    #1%
    \ifbool{cbx:parens}
      {\bibcloseparen\global\boolfalse{cbx:parens}}
      {}}}
  \savebibmacro{cite}
  \savebibmacro{textcite}
  \renewbibmacro*{cite}{%
    \printtext[citehyperref]{%
      \restorebibmacro{cite}%
      \usebibmacro{cite}}}
  \renewbibmacro*{textcite}{%
    \ifboolexpr{
      ( not test {\iffieldundef{prenote}} and
        test {\ifnumequal{\value{citecount}}{1}} )
      or
      ( not test {\iffieldundef{postnote}} and
        test {\ifnumequal{\value{citecount}}{\value{citetotal}}} )
    }
      {\DeclareFieldAlias{textcitehyperref}{noformat}}
      {}%
    \printtext[textcitehyperref]{%
      \restorebibmacro{textcite}%
      \usebibmacro{textcite}}}

  % Remove editor form everything but books
  \AtEveryBibitem{%
   \ifentrytype{book}{}{%
    \clearname{editor}
   }
  }
}{
  %-----------------------------------------------------------------------------
  %  If using bibtex
  %-----------------------------------------------------------------------------
  % Load required packages and select citation style
  \usepackage[round]{natbib}
  \bibliographystyle{apalike}
}

%------------------------------------------------------------------------------
%	Some Styling
%------------------------------------------------------------------------------
% Link colors
\definecolor{linkcolor}{HTML}{000000}
\hypersetup{
  colorlinks = true,
  linkcolor  = linkcolor,
  urlcolor   = linkcolor,
  citecolor  = linkcolor,
}

% Changing the style of captions in figures etc.
\captionsetup{labelfont=bf, format=plain, font=small}

% Change how equations are referenced
\renewcommand{\theequation}{Equation \arabic{equation}}%

% Avoid skip when including text from external .tex file
\newcommand{\inputy}[1]{\input{#1}\unskip}

% For supplementary material
\newcommand{\beginappendix}{%
  \setcounter{table}{0}
  \renewcommand{\thetable}{S\arabic{table}}%
  \setcounter{figure}{0}
  \renewcommand{\thefigure}{S\arabic{figure}}%
  \setcounter{equation}{0}
  \renewcommand{\theequation}{Equation S\arabic{equation}}%
  \setcounter{section}{0}
  \renewcommand{\thesection}{A.\arabic{section}}%
}

%------------------------------------------------------------------------------
%  ORCID
%------------------------------------------------------------------------------
% Creating some TikZ styles
\tikzset{
  nonterminal/.style = {rectangle
    , minimum size = 6mm
    , very thick
    , draw = black!
  }
}

% Create the ORCID logo
\definecolor{orcidlogocol}{HTML}{A6CE39}
\tikzset{
  orcidlogo/.pic={
    \fill[orcidlogocol]
      svg{M256,128c0,70.7-57.3,128-128,128C57.3,256,0,198.7,0,128C0,57.3,57.3
        ,0,128,0C198.7,0,256,57.3,256,128z};
    \fill[white]
      svg{M86.3,186.2H70.9V79.1h15.4v48.4V186.2z}
      svg{M108.9,79.1h41.6c39.6,0,57,28.3,57,53.6c0,27.5-21.5,53.6-56.8
        ,53.6h-41.8V79.1z M124.3,172.4h24.5c34.9,0,42.9-26.5
        ,42.9-39.7c0-21.5-13.7-39.7-43.7-39.7h-23.7V172.4z}
      svg{M88.7,56.8c0,5.5-4.5,10.1-10.1,10.1c-5.6,0-10.1-4.6-10.1-10.1c0-5.6
      ,4.5-10.1,10.1-10.1C84.2,46.7,88.7,51.3,88.7,56.8z};
  }
}

% Command to create the ORCID
\newcommand\orcid[1]{\href{https://orcid.org/#1}{\mbox{\scalerel*{
\begin{tikzpicture}[yscale=-1,transform shape]
  \pic{orcidlogo};
\end{tikzpicture}
}{|}}}}

%------------------------------------------------------------------------------
%	Titlepage: Header
%------------------------------------------------------------------------------
\title{\textbf{Appendix}\\ Methods for Implementing Integrated Step-Selection
Functions with Incomplete Data}

% List of Authors
\author{
  David D. Hofmann\textsuperscript{1,2,\S} \orcid{0000-0003-3477-4365} \and
  Gabriele Cozzi\textsuperscript{1,2} \orcid{0000-0002-1744-1940} \and
  John Fieberg \textsuperscript{3} \orcid{0000-0002-3180-7021}
}

% Reduce spacing between authors
\makeatletter
\def\and{%
  \end{tabular}%
  \hskip -0.5em \@plus.17fil\relax
  \begin{tabular}[t]{c}}
\makeatother

% Current Date
% \date{\today}

% And here the masterpiece begins
\begin{document}

% Change page numbering
\pagenumbering{gobble}

% Create Titlepage
\maketitle

%------------------------------------------------------------------------------
%	Titlepage: Additional Info
%------------------------------------------------------------------------------
\begin{flushleft}

\vspace{0.5cm}

\textsuperscript{1} Department of Evolutionary Biology and Environmental
Studies, University of Zurich, Winterthurerstrasse 190, 8057 Zurich,
Switzerland.

\textsuperscript{2} Botswana Predator Conservation Program, Wild Entrust,
Private Bag 13, Maun, Botswana.

\textsuperscript{3} Department of Fisheries, Wildlife, and Conservation Biology,
University of Minnesota, St. Paul, MN, USA.

\textsuperscript{\S} Corresponding author: \href{mailto://david.hofmann2@uzh.ch}{david.hofmann2@uzh.ch}

\vspace{4cm}

\textbf{Running Title:} Step-Selection Analyses with Missing Data

\vspace{0.5cm}

\textbf{Keywords:} animal movement, gps data, imputation, incomplete data,
missing fixes, step-selection analyses, step-selection functions

\end{flushleft}

%------------------------------------------------------------------------------
%	Main Text
%------------------------------------------------------------------------------
\newpage

% Change page numbering
\pagenumbering{arabic}

% Create linenumbers
% \linenumbers

% Change to appendix counters
\appendix
\beginappendix

%------------------------------------------------------------------------------
%	Appendix S1:
%------------------------------------------------------------------------------
\newpage
\section{Landscape Simulation: Different Autocorrelation Scenarios}
\begin{figure}[!ht]
  \begin{center}
  \includegraphics[width = \textwidth]{99_CovariatesAppendix.png}
  \caption{Simulated landscapes under different levels of autocorrelation (5,
  10, 20; from top to bottom). Autocorrelation only affected the layers
  \textsf{elev} and \textsf{forest}, which were both simulated using a Gaussian
  random field neutral landscape model \citep{Schlather.2015} using the
  \texttt{R}-package \texttt{NLMR} \citep{Sciaini.2018}. Simulations were
  repeated 100 times for each autocorrelation scenario, thus resulting in 300
  unique landscape configurations.}
  \label{Covariates}
  \end{center}
\end{figure}

%------------------------------------------------------------------------------
%	Appendix S2:
%------------------------------------------------------------------------------
\newpage
\section{Dynamic Tentative Distribution Parameters}
\begin{figure}[!ht]
  \begin{center}
  \includegraphics[width = 0.9\textwidth]{99_DistributionParameters.png}
  \caption{Tentative prameter estimates for the von Mises distribution (top row) and
  gamma distribution (bottom row) fitted to steps with different durations. The
  von Mises distribution requires one parameter, namely a concentration
  parameter ($\kappa$). The gamma distribution requires a shape parameter ($k$)
  and a scale parameter ($\theta$). The subscript $_0$ is used to indicate that
  these are tentative distribution parameters (sensu \citealp{Avgar.2016} and
  \citealp{Fieberg.2021}).}
  \label{DistributionParameters}
  \end{center}
\end{figure}

%------------------------------------------------------------------------------
%	Appendix S3:
%------------------------------------------------------------------------------
\newpage
\section{Distribution of Relative Turning Angles following Different Step Durations}
\begin{figure}[!ht]
  \begin{center}
  \includegraphics[width = 0.9\textwidth]{99_TurningAnglePreviousDuration.png}
  \caption{Density of relative turning angles associated with steps of $\Delta t
  = 1$ following steps with different durations for all three autocorrelation
  scenarios (5, 10, and 20). To generate this figure, we assumed a missingness
  of 0.5 and forgiveness of 5.}
  \label{PreviousDuration}
  \end{center}
\end{figure}

%------------------------------------------------------------------------------
%	Appendix S4:
%------------------------------------------------------------------------------
\newpage
\section{Model Estimates across all Scenarios}
\begin{figure}[!ht]
  \begin{center}
  \includegraphics[width = \textwidth]{99_ModelComparisonAppendix.png}
  \caption{Parameter estimates across different autocorrelation scenarios (5,
  10, 20; panels a, b, and c) and missingness levels (0\% - 50\%; from left
  to right). True simulation parameters are indicated by the solid black
  lines. Parameter estimates from the different approaches are given by the
  colored symbols, and their bootstrap 95\% CIs across 100 replicates
  by the colored lines.}
  \label{ModelComparison}
  \end{center}
\end{figure}

%------------------------------------------------------------------------------
%	Appendix S5:
%------------------------------------------------------------------------------
\newpage
\section{Case Study Covariates}

\begin{figure}[htpb]
  \begin{center}
  \includegraphics[width = \textwidth]{99_CaseStudyCovariates.png}
  \caption{Covariates used for the case study, overlaid with the GPS data of a
  spotted hyena called ``Apollo'' (lines in black). Apollo was originally
  collared in 2007 in northern Botswana and monitored until 2011. The depicted
  area is part of the Okavango Delta, which is a massive wetland area. Data was
  projected to a local projection (EPSG:32734).}
  \label{CaseStudyCovariates}
  \end{center}
\end{figure}

%------------------------------------------------------------------------------
%	Appendix S6:
%------------------------------------------------------------------------------
\newpage
\section{Case Study Model Output}

\begin{table}[htpb]
  \begin{center}
  \caption{Model results from the case study using GPS data collected on Apollo.
  In F1, forgiveness was set to one (only 2-hour steps were considered), whereas
  in F3-S and F3-SH a forgiveness of three was employed (allowing for step
  durations of up to 6 hours). In model F3-S, the step duration was interacted
  with step descriptors. In model F3-SH, step duration was interacted with step
  descriptors and habitat covariates.}
  \label{CaseStudyTable}
    \begin{adjustbox}{max width=0.70\textwidth}
    \begin{threeparttable}[h]
      \input{99_CaseStudyTable.tex}
      \begin{tablenotes}
       \item \textit{Significance codes: * \(p < 0.10\) \quad ** \(p < 0.05\)
        \quad *** \(p < 0.01\)}
      \end{tablenotes}
    \end{threeparttable}
  \end{adjustbox}
  \end{center}
\end{table}

\newpage
\begingroup
\singlespacing
\ifthenelse{\boolean{usebiblatex}}{
  \begin{refcontext}[sorting=nyt]
  \printbibliography
  \end{refcontext}
}{
  \bibliography{LiteratureBibtex}
}
\endgroup
\end{document}
